\title{xLSTM: Extended Long Short-Term Memory}

\begin{abstract}
The introduction of the constant error carousel and gating mechanisms during the 1990s laid the foundation for Long Short-Term Memory (LSTM) networks. Over the years, LSTMs have demonstrated considerable resilience and have powered a wide range of breakthroughs in deep learning, notably serving as the backbone of the earliest Large Language Models (LLMs). Nevertheless, with the emergence of Transformer architectures—characterized by highly parallelizable self-attention—there has been a paradigm shift, resulting in the eclipse of LSTMs at scale. This work revisits a foundational question: What levels of performance are achievable in language modeling by scaling LSTMs to billions of parameters, applying advanced methodologies from contemporary LLMs while systematically addressing historical LSTM constraints? To this end, we introduce exponential gating, paired with normalization and stabilization techniques suited to large-scale training. Additionally, we re-envision the LSTM memory construct, presenting: (i) sLSTM, which utilizes a scalar-valued memory, scalar updates, and an updated memory mixing method, and (ii) mLSTM, a design featuring a matrix memory and a covariance-based update, enabling full parallelization. By embedding these LSTM variants into the structure of residual blocks, we assemble xLSTM blocks, which are subsequently composed in deep residual stacks, forming xLSTM architectures. The combination of exponential gating and restructured memory substantially enhances xLSTM models, allowing them to match or exceed the performance and scaling behavior of leading Transformers and State Space Models.\\
Code available at: \url{https://github.com/NX-AI/xlstm}
\end{abstract}

\section{Introduction}
\label{sec:intro}

The Long Short-Term Memory (LSTM) framework~\citep{Hochreiter:91,Hochreiter:97nips,Hochreiter:97}, featuring the constant error carousel and gate mechanisms, was developed specifically to address the vanishing gradient issue observed in recurrent neural networks~\citep{Hochreiter:91,Hochreiter:00book}:

\begin{align}
\label{eq:lstm_idea}
  c_t \ = \ f_t \ c_{t-1} \ + \ 
          i_t \  z_t \ , \quad  
          h_t \ = \ o_t \ \psi({c_t}) \ . 
\end{align}

The cell state $c_{t-1}$ (represented in green) is updated additively with cell inputs $z_t$ in the mechanism known as the constant error carousel, with the process being regulated by sigmoid gates (depicted in blue). The forget gate $f_t$ and input gate $i_t$ are responsible for determining how this update proceeds, whereas the output gate $o_t$ dictates the release of information from the memory cell, corresponding to the hidden state $h_t$. The function $\psi$ acts to squash or normalize the cell state prior to the output gating, which ultimately produces the hidden state.

LSTMs have found widespread success across numerous fields \citep{Hochreiter:01,Hochreiter:07,Schmidhuber:15}, maintaining dominance in text generation tasks until the introduction of Transformers in 2017~\citep{Vaswani:17}. Their capability has been illustrated in a variety of sequence-centric problems, including but not limited to, text synthesis~\citep{Graves:13,Karpathy:15blog}, handwriting production~\citep{Graves:13}, sequence-to-sequence translation~\citep{Sutskever:14nips}, code evaluation~\citep{Zaremba:14arxiv}, image caption creation~\citep{Karpathy:15,Hossain:19}, automated source code generation~\citep{Karpathy:15blog}, rainfall-runoff prediction~\citep{Kratzert:18,Kratzert:19}, and hydrological modeling for flood alerts~\citep{Nearing:24}. Within reinforcement learning, LSTMs continue to serve as high-performing sequence models, as evidenced by their role in systems such as AlphaStar for StarCraft II~\citep{Vinyals:17short}, OpenAI Five for Dota 2~\citep{OpenAI:19}, and magnetic controller models for nuclear fusion~\citep{Degrave:22short}. A particular strength of LSTMs lies in their ability to learn and abstract high-level semantic features, efficiently storing such representations within their memory cells~\citep{Karpathy:15blog}, an ability highlighted through findings such as number and syntax units~\citep{Lakretz:19}, linguistic neurons~\citep{Bau:19}, and sentiment neurons~\citep{Radford:17arxiv}. LSTMs remain integral to critical contemporary systems~\citep{Degrave:22short,Nearing:24}, demonstrating enduring utility and robustness over time.

Although LSTMs have achieved significant breakthroughs, they exhibit three core shortcomings: (i) Lack of capacity to revise stored information. We demonstrate this shortcoming using the {\it Nearest Neighbor Search} task (see Appendix~\ref{sec:appNearestNeighborSearch} for further details): given a reference vector, the model must sequentially inspect a series of vectors to retrieve the associated value of the most similar one at the end of the sequence. The left panel of Figure~\ref{fig:lstmProblems} presents the mean squared error results for this challenge. LSTM networks exhibit difficulty overwriting previously stored values when subsequently encountering a more similar vector. In contrast, our proposed xLSTM, through exponential gating, overcomes this deficit. (ii) Restriction in storage capacity, as conventional LSTMs must represent information within scalar-valued cell states. We illustrate this with the {\it Rare Token Prediction} task. The right panel of Figure~\ref{fig:lstmProblems} reports token prediction perplexity on Wikitext-103~\citep{Merity:17}, broken down by token frequency groups. Due to their constrained memory capacity, LSTMs yield higher perplexity, especially on infrequent tokens. The xLSTM, leveraging a matrix-based memory, resolves this issue. (iii) Intrinsic limits on parallel computation stemming from recurrent memory mixing, that is, the dependence between hidden states at consecutive time steps forces strict sequential processing.

The shortcomings of LSTM have led to the development and widespread adoption of Transformers~\citep{Vaswani:17} for language modeling. An open question is how language modeling performance would be affected if these LSTM constraints were addressed and LSTM architectures were scaled to match the magnitude of contemporary Large Language Models.

\section{Extended Long Short-Term Memory}
\label{sec:xLSTM}

In order to address the constraints inherent in standard LSTM architectures, Extended Long Short-Term Memory (xLSTM) proposes two fundamental enhancements to the LSTM conceptual framework described in Equation~\eqref{eq:lstm_idea}. These enhancements—namely, exponential gating mechanisms and innovative memory configurations—result in two new LSTM variants: (i) the sLSTM (detailed in Section~\ref{sec:sLSTM}), which employs scalar-valued memory, scalar updates, and incorporates a mechanism for memory mixing; and (ii) the mLSTM (see Section~\ref{sec:mLSTM}), which operates with matrix memory structures and utilizes a covariance-based (outer product) memory update that supports full parallelization. Both sLSTM and mLSTM introduce exponential gating to improve upon traditional LSTM dynamics. In pursuit of parallel computation, the mLSTM eliminates the use of memory mixing, specifically by removing recurrent connections among hidden units. Both mLSTM and sLSTM architectures may be generalized to accommodate multiple memory cells, wherein sLSTM exhibits memory mixing across different cells. Additionally, sLSTM can be equipped with multiple heads; in this scenario, memory mixing does not occur across heads but is instead restricted to cells within the same head. This framework—combining multi-head configuration and exponential gating—introduces a novel form of memory mixing for sLSTM. Conversely, for mLSTM, having multiple heads is functionally identical to having multiple memory cells.

By embedding these novel LSTM variants within residual block modules, we obtain what are termed xLSTM blocks (refer to Section~\ref{sec:xLSTMblock}). When these xLSTM blocks are stacked residually within a network, the resulting models are referred to as xLSTM architectures (refer to Section~\ref{sec:xLSTMarchitecure}). The structure and components of the xLSTM architecture are illustrated in Figure~\ref{fig:xlstm_sketch}.

\subsection{Review of the Long Short-Term Memory}
\label{sec:orgLSTM}

The foundational concept behind LSTM~\citep{Hochreiter:91,Hochreiter:97nips,Hochreiter:97} centers around a scalar memory cell, which functions as a key processing and storage unit designed to counteract the problem of vanishing gradients~\citep{Hochreiter:91,Hochreiter:00book} via the constant error carousel (cell state update mechanism). The architecture of this memory cell is governed by three types of gates: the input gate, the output gate, and the forget gate, the latter of which was incorporated by \citet{Gers:00}. The dynamic behavior of the LSTM memory cell at each time step $t$ can be described by the following update equations:

\begin{align}
c_t \ &= \  f_t \ c_{t-1} \ + \ i_t \ z_t &  & &\text{cell state} \\
h_t  \ &= \ o_t \ \tilde{h}_t \ , & \tilde{h}_t \ &= \ \psi \left( c_t\right)
  &\text{hidden state} \\
z_t \ &= \ \varphi \left( \tilde{z}_t \right) \ , 
  &\tilde{z}_t \ &=  \ \mathbf{w}^\top_{z} \ \mathbf{x}_t \ + \
  r_{z}  h_{t-1} \ + \  b_{z} \ \
  &\text{cell input} \\
i_t \ &= \ \sigma \left( \tilde{i}_t  \right) \ , 
  &\tilde{i}_t \ &= \ \mathbf{w}^\top_{i} \ \mathbf{x}_t \ + \
  r_{i}  \ h_{t-1} \ + \  b_{i} \ \
  &\text{input gate} \\
f_t \ &= \ \sigma \left(  \tilde{f}_t \right) \ , 
  &\tilde{f}_t \ &= \ \mathbf{w}^\top_{f} \ \mathbf{x}_t  \ + \
  r_{f}  \ h_{t-1} \ + \  b_{f} \ \
  &\text{forget gate} \\
o_t \ &= \ \sigma \left( \tilde{o}_t \right) \ , 
  &\tilde{o}_t  \ &= \ \mathbf{w}^\top_{o} \ \mathbf{x}_t \ + \
  r_{o}  \ h_{t-1} \ + \  b_{o} \ \
  &\text{output gate} 
\end{align}

The vectors $\mathbf{w}_{z}$, $\mathbf{w}_{i}$, $\mathbf{w}_{f}$, and $\mathbf{w}_{o}$ denote the input weight vectors that link the inputs $\mathbf{x}_t$ to the cell input, input gate, forget gate, and output gate, respectively. The parameters $r_{z}$, $r_{i}$, $r_{f}$, and $r_{o}$ represent the recurrent weights that connect the previous hidden state $h_{t-1}$ to the cell input, input gate, forget gate, and output gate, respectively. The terms $b_{z}$, $b_{i}$, $b_{f}$, and $b_{o}$ serve as the bias components for each corresponding gate or input. The functions $\varphi$ and $\psi$ refer to the nonlinearities applied to the cell input and the hidden state (commonly $\tanh$). The role of $\psi$ is to constrain or scale the cell state, as it can potentially diverge without such a nonlinearity. The gate activations all utilize the sigmoid function, given by $\sigma \left( x \right)= 1/(1+\exp(-x))$. In subsequent versions, memory cells are organized as vectors $\mathbf{c}_t \in \mathbb{R}^{d}$, making it possible to introduce recurrent weight matrices $\mathbf{R} \in \mathbb{R}^{d \times d}$ to combine and transform outputs across memory cells~\citep{Greff:15}; additional information is provided in Appendix~\ref{sec:apporgLSTM}. Studies removing individual components have demonstrated that each part of the memory cell architecture plays an essential role~\citep{Greff:15}.

\subsection{sLSTM}
\label{sec:sLSTM}

To enable LSTMs to update their memory selectively, we propose the integration of exponential gates (illustrated in red), combined with mechanisms for normalization and stabilization. Specifically, both the input and forget gates are allowed to utilize exponential activation functions. For the normalization process, we define a normalizer state, which accumulates the results of multiplying the input gate by all subsequent forget gates. 

The forward pass for the scalar sLSTM is defined as follows:

\begin{align}
\label{eq:slstmforward}
c_t \ &= \  f_t \ c_{t-1} \ + \ i_t \ z_t & & 
 &\text{cell state} \\
n_t \ &= \  f_t \ n_{t-1} \ + \ i_t  & & 
 &\text{normalizer state} \\
h_t \ &= \ o_t \ \tilde{h}_t \ ,
  &\tilde{h}_t \ &= \ c_t / n_t
&\text{hidden state} \\
z_t \ &= \ \varphi \left( \tilde{z}_t \right) \ , 
  &\tilde{z}_t \ &=  \ \mathbf{w}^\top_{z} \ \mathbf{x}_t \ + \
  r_{z}  h_{t-1} \ + \  b_{z}
  &\text{cell input} \\
i_t \ &= \ \!\exp \left( \tilde{i}_t  \right) \ , 
  &\tilde{i}_t \ &= \ \mathbf{w}^\top_{i} \ \mathbf{x}_t \ + \
  r_{i}  \ h_{t-1} \ + \  b_{i}
  &\text{input gate} \\
f_t \ &= \ \sigma \left(  \tilde{f}_t \right) \ \text{OR} \
\!\exp \left(  \tilde{f}_t \right) \ ,
  &\tilde{f}_t \ &= \ \mathbf{w}^\top_{f} \ \mathbf{x}_t  \ + \
  r_{f}  \ h_{t-1} \ + \  b_{f}
  &\text{forget gate} \\
o_t \ &= \ \sigma \left( \tilde{o}_t \right) \ , 
  &\tilde{o}_t  \ &= \ \mathbf{w}^\top_{o} \ \mathbf{x}_t \ + \
  r_{o}  \ h_{t-1} \ + \  b_{o}
  &\text{output gate} 
\end{align}

The original LSTM gating mechanisms, which involve input- and/or hidden-state-dependent gates along with a bias term, are adapted for these novel architectures. Since exponential activation functions have the potential to generate extremely large outputs, resulting in overflow issues, we introduce an extra state \mbox{$m_t$~\citep{Milakov:18arxiv}} to stabilize the gates:

\begin{align}
\label{eq:slstmstabil}
m_t \ &= \ \max \left( \log ( f_t ) + m_{t-1} , \log ( i_t ) \right) &\text{stabilizer state} \\
i'_t \ &= \ \!\exp \left( \log \left ( i_t \right) - m_t \right) \ = \ \!\exp \left( \tilde{i}_t - m_t \right) \ \, 
  &\text{stabil. input gate} \\
f'_t \ &= \ \!\exp \left( \log \left( f_t \right) + m_{t-1} - m_t \right) \ \, 
  &\text{stabil. forget gate}
\end{align}

As demonstrated in Appendix~\ref{sec:appsLSTM}, substituting $f_t$ with $f'_t$ and $i_t$ with $i'_t$ during the forward pass leaves both the network's output and the loss gradients with respect to the parameters unchanged.

\paragraph{New Memory Mixing.} Similar to standard LSTM architectures, sLSTM supports the use of several memory cells (refer to Appendix~\ref{sec:appsLSTM}). The presence of multiple memory cells facilitates memory mixing through the recurrent connections $\mathbf{R}_{\mathbf{z}}$, $\mathbf{R}_{\mathbf{i}}$, $\mathbf{R}_{\mathbf{f}}$, $\mathbf{R}_{\mathbf{o}}$, which transmit information from the hidden state vector $\mathbf{h}$ to the memory cell input $\mathbf{z}$ as well as to the gates $\mathbf{i}$, $\mathbf{f}$, and $\mathbf{o}$. A distinguishing characteristic of the new memory mixing approach is the application of exponential gating. In the sLSTM, each head can implement memory mixing internally, but not between different heads. Thus, by integrating head structures in sLSTM in conjunction with exponential gating, a novel paradigm for memory mixing is realized.

\subsection{mLSTM}
\label{sec:mLSTM}

In order to improve the memory capabilities of LSTMs, we transition from representing the LSTM cell as a scalar value $c \in \mathbb{R}$ to expressing it as a matrix $\mathbf{C} \in \mathbb{R}^{d \times d}$. As a result, the retrieval process utilizes matrix multiplication. At each time step $t$, the model stores two vectors: a key $\mathbf{k}_t \in \mathbb{R}^d$ and a value $\mathbf{v}_t \in \mathbb{R}^d$, following the terminology used in Transformer architectures. Subsequently, at time $t + \tau$, the stored value $\mathbf{v}_t$ can be accessed using a query vector $\mathbf{q}_{t+\tau} \in \mathbb{R}^d$. This paradigm aligns with the framework of Bidirectional Associative Memories (BAMs) as described by \citep{Kohonen:72,Anderson:72,Nakano:72,Anderson:77}.

The method for updating the covariance to encode a key-value pair is given by~\citep{Sejnowski:77,Dayan:91}

\begin{align}
\mathbf{C}_t \ &= \ \mathbf{C}_{t-1} \ + \ \mathbf{v}_{t} \ \mathbf{k}_t^\top \ .
\end{align}

We consider a scenario where layer normalization is applied prior to projecting the inputs into keys and values, ensuring that these projections are mean-centered at zero. The rule for updating the covariance is provably optimal~\citep{Dayan:91} with respect to maximizing the separability between retrieved binary vectors, which is tantamount to achieving the highest possible signal-to-noise ratio. Enhanced separability can be achieved if retrieval is confined to pairwise associations and if one accepts the quadratic computational cost characteristic of attention mechanisms~\citep{Krotov:16,Krotov:17,Ramsauer:21}. The covariance update mechanism also coincides with the formulation used in Fast Weight Programmers~\citep{Schmidhuber:92ncfastweights,Schlag:21}, which were later modified to include a fixed decay coefficient applied to $\mathbf{C}_{t-1}$ and a fixed learning coefficient applied to $\mathbf{v}_{t} \mathbf{k}_t ^\top$~\citep{Ba:16}. In line with this reasoning, we incorporate the covariance update scheme within the LSTM architecture, mapping the forget gate to the decay coefficient, the input gate to the learning coefficient, and using the output gate to modulate the retrieved vector.

Within this matrix memory framework, the normalizer state represents a sum of key vectors, each scaled according to the product of the input gate and subsequent forget gates. This state effectively tracks the accumulated influence of the gates over time. Because the dot product between a query and the normalizer state may sometimes approach zero, its absolute value is taken and constrained to be at least a minimum threshold (commonly 1.0), as described in earlier work~\citep{Sun:23arxiv}. The forward computation for mLSTM is as follows:

\begin{align}
\label{eq:mlstm_recurrent_begin}
\mathbf{C}_t \ &= \  f_t \ \mathbf{C}_{t-1} \ + \ 
  i_t \ \mathbf{v}_t \ \mathbf{k}_t^\top & &  &\text{cell state} \\
\mathbf{n}_t \ &= \  f_t \ \mathbf{n}_{t-1} \ + \ 
  i_t \ \mathbf{k}_t & &  &\text{normalizer state} \\
\mathbf{h}_t  \ &= \ \mathbf{o}_t \ \odot \ \tilde{\mathbf{h}}_t
\ , \qquad \qquad 
\tilde{\mathbf{h}}_t \ = \   \mathbf{C}_t \mathbf{q}_t \ / \ 
  \max \left\{ |\mathbf{n}_t^\top \mathbf{q}_t|, 1 \right\} 
   &&&\text{hidden state} \\
\mathbf{q}_t \ &= \ \mathbf{W}_q \ \mathbf{x}_t \ + \ \mathbf{b}_q  & &  &\text{query input} \\
\mathbf{k}_t \ &= \ \frac{1}{\sqrt{d}} \mathbf{W}_k \ \mathbf{x}_t \ + \ \mathbf{b}_k  & &  &\text{key input} \\
\mathbf{v}_t \ &= \ \mathbf{W}_v \ \mathbf{x}_t \ + \ \mathbf{b}_v  & &  &\text{value input} \\
i_t \ &= \ \!\exp \!  \! \left( \tilde{i}_t  \right) \ , 
 \qquad \qquad \ \ \ \ \,
  \tilde{i}_t \ = \ \mathbf{w}^\top_{i} \ \mathbf{x}_t \ + \  b_{i} \ \ \ 
  &&&\text{input gate} \\
f_t \ &= \ \sigma \!  \left(  \tilde{f}_t \right) \ \text{OR} \ \!\exp \!  \! \left(  \tilde{f}_t \right) , \ \ \: \!
\tilde{f}_t \ = \ \mathbf{w}^\top_{f} \ \mathbf{x}_t  \ + \
  b_{f}
  &&& \text{forget gate} \\
\label{eq:mlstm_recurrent_end}
\mathbf{o}_t \ &= \ \sigma \left( \tilde{\mathbf{o}}_t \right) \ , \qquad \qquad \quad \ \ \ \,
  \tilde{\mathbf{o}}_t  \ = \ \mathbf{W}_{\mathbf{o}} \ \mathbf{x}_t \ + \
  \mathbf{b}_{\mathbf{o}}
  &&&\text{output gate} 
\end{align}

The architecture of mLSTM permits several memory cells, similar to the standard LSTM. In the context of mLSTM, employing multiple heads or multiple cells is effectively identical, due to the absence of memory mixing between them. To ensure stable behavior of mLSTM's exponential gates, we apply the stabilization methods utilized for sLSTM (refer to Equation~\ref{eq:slstmstabil}). Given that memory mixing does not occur in mLSTM, its recurrence can be expressed in a parallelized format. Further information is provided in Appendix~\ref{sec:appmLSTM}.

\subsection{xLSTM Architecture}
\label{sec:xLSTMarch}

\paragraph{xLSTM Blocks.}
\label{sec:xLSTMblock}
The xLSTM block is designed to generate a non-linear summary of preceding information within a high-dimensional representation, enhancing the model's ability to distinguish between varying histories or contextual settings. Effective disambiguation of past contexts is essential for accurately forecasting subsequent elements in a sequence, such as the next token. We draw upon Cover’s Theorem~\citep{Cover:65}, which asserts that patterns mapped non-linearly into higher-dimensional spaces are more apt to become linearly separable compared to their arrangement in the initial, lower-dimensional domain. We investigate two types of residual block configurations: (i) The residual block with a post up-projection (mirroring Transformer architectures), where a non-linear summary is constructed in the original dimension, followed by a linear projection into a space of higher dimensionality, then a non-linear activation is introduced, and finally a linear mapping restores the original dimensionality; refer to the leftmost diagram in Figure~\ref{fig:Backbones} and the third column in Figure~\ref{fig:xlstm_sketch}. A comprehensive schematic is available in Figure~\ref{fig:appsLSTM_detailed} in the appendix. (ii) The residual block with a pre up-projection (akin to State Space Models), in which the data are first linearly projected to a space with greater dimensionality,

The past is first non-linearly encoded within a high-dimensional space before being linearly projected back into the original space. In the context of an xLSTM block utilizing an sLSTM, we employ the post up-projection block. Conversely, when an xLSTM block incorporates an mLSTM, the pre up-projection block is chosen, as the high-dimensional representation enables greater memory capacity. Further specifics can be found in the left panel of Figure~\ref{fig:Backbones}, the third column of Figure~\ref{fig:xlstm_sketch}, and in Figure~\ref{fig:appsLSTM_detailed} within the appendix.

\paragraph{xLSTM Architecture. }
\label{sec:xLSTMarchitecure}
The xLSTM model is formed by assembling its foundational units in a residual stacking manner~\citep{Srivastava:15,He:16}. Our approach employs the widely adopted pre-LayerNorm~\citep{Ba:16b} residual backbone, mirroring the standard practice in modern Large Language Models. For a visual representation, refer to the last column of Figure~\ref{fig:xlstm_sketch}.

\subsection{Memory and Speed Considerations}

Unlike Transformers, xLSTM networks exhibit linear computational complexity and maintain constant memory usage regardless of sequence length. The compressive nature of xLSTM memory makes it particularly advantageous for deployment in industrial contexts and for edge computing scenarios.

Although the mLSTM's memory component is parameter-free, its reliance on a $d \times d$ matrix for both storage and updates results in substantial computational demands. We balance the richness of the memory capacity with the associated increase in computational requirements. However, as these operations can be parallelized on GPUs, their impact on actual processing time remains minimal.

Although mLSTM can be parallelized similarly to FlashAttention~\citep{Dao:22,Dao:23} or GLA~\citep{Yang:23arxiv}, the architecture of sLSTM prevents parallel execution because of the memory mixing inherent in its hidden-hidden connections. Nonetheless, we have created an efficient CUDA implementation that leverages GPU memory optimization down to the register level, resulting in performance that is generally less than twice as slow as that of mLSTM.

\section{Related Work}
\label{sec:related}

\paragraph{Linear Attention.} Multiple approaches have been proposed to address the Transformer's quadratic computational burden with respect to context length, aiming to reduce attention computation to linear complexity. The Synthesizer method dispenses with token-to-token interactions, instead learning artificial attention weights~\citep{Tay:20arxiv}. In Linformer, self-attention is approximated by projecting inputs into a low-rank space, resulting in a linear-time computation~\citep{Wang:20arxiv}. Linear Transformer achieves a linear formulation of the attention operation~\citep{Katharopoulos:20}. The Performer approximates the attention softmax function linearly using positive orthogonal random feature techniques~\citep{Choromanski:21}. Additionally, alternatives to traditional attention, such as fast long convolutional operations, are implemented in Structured Global Convolution (SGConv)~\citep{Li:22} and in the Hyena Hierarchy~\citep{Poli:23}.

\paragraph{State Space Models.} State Space Models (SSMs) have garnered considerable attention in recent years due to their linear computational complexity with respect to sequence length and their competitive results against Transformer architectures. Early notable models include the Structured State Space sequence model (S4)~\citep{Gu:21}. Subsequent advancements introduced the Diagonal State Space (DSS) model~\citep{Gupta:22}, Gated State Space (GSS) models~\citep{Mehta:22}, S5 model~\citep{Smith:22}, Bidirectional Gated SSM (BiGS)~\citep{Wang:22}, H3 model~\citep{Fu:23}, and most recently, Mamba~\citep{Gu:24arxiv}.

\paragraph{Recurrent Neural Networks.} In recent research, Recurrent Neural Networks (RNNs) have been proposed as alternatives to Transformers and attention-based mechanisms, owing to their computational efficiency that scales linearly with context length. Variants such as Deep Linear Recurrent Units (LRUs) have demonstrated strong language modeling performance~\citep{Orvieto:23, De:24arxiv}, alongside models like Hierarchically Gated Linear RNN (HGRN)~\citep{Qin:23} and its successor HGRN2~\citep{Qin:24arxiv}. Among RNN-based models for large-scale language modeling, RWKV~\citep{Peng:23arxivshort,Peng:24arxivshort} has gained prominence, achieving results comparable to Transformer architectures.

\paragraph{Gating.}
A fundamental concept underlying LSTM is gating, which has since been rediscovered and reframed in numerous recent works. Variants of gating have been incorporated into models such as HGRN~\citep{Qin:23}, HGRN2~\citep{Qin:24arxiv}, Gated Linear Attention (GLA)~\citep{Yang:23arxiv}, Gated State Space (GSS) models~\citep{Mehta:22}, Bidirectional Gated SSM (BiGS)~\citep{Wang:22}, Moving Average Equipped Gated Attention (MEGA)~\citep{Ma:22}, RWKV~\citep{Peng:23arxivshort}, and Mamba~\citep{Gu:24arxiv}.

\paragraph{Covariance Update Rule.} In order to improve storage capacity, we incorporated a matrix memory into the mLSTM cell utilizing a covariance update rule. Similar update strategies have been employed by models such as Fast Weight Programmers \citep{Schmidhuber:92ncfastweights,Schlag:21}, RWKV-5 and RWKV-6~\citep{Peng:24arxivshort}, Retention~\citep{Sun:23arxiv}, Linear Transformer~\citep{Katharopoulos:20}, and HGRN2~\citep{Qin:24arxiv}.

\paragraph{Most Related.} The models that are most closely aligned with xLSTM from a conceptual standpoint include Retention~\citep{Sun:23arxiv}, RWKV~\citep{Peng:23arxivshort,Peng:24arxivshort}, and HGRN2~\citep{Qin:24arxiv}. These frameworks incorporate elements such as matrix memory and gating mechanisms. Nevertheless, unlike the newly introduced sLSTM, these models lack the ability to perform memory mixing. This capability is essential for effectively handling state tracking tasks, which renders LSTMs more powerful than State Space Models (SSMs) and Transformers~\citep{Merrill:24,Deletang:23} in such scenarios. State tracking plays a critical role in tasks such as code execution and maintaining continuity of entities over extended narratives.

\paragraph{Residually Stacking Architectures.} Consistent with the design principles of nearly all modern large-scale deep learning models, xLSTM architectures are formed by stacking modules in a residual manner~\citep{Srivastava:15,He:16}.

This strategy of residual stacking has played a crucial role in the development of deep convolutional neural networks~\citep{He:16} as well as Transformers~\citep{Vaswani:17}. Transformers, in particular, underpin the capabilities of Large Language Models (LLMs), such as GPT-3~\citep{Brown:20short}, ChatGPT~\citep{Schulman:22short}, GPT-4~\citep{Achiam:23gpt4short}, Megatron-LM~\citep{Shoeybi:19}, Gopher~\citep{Rae:21short}, ERNIE 3.0 Titan~\citep{Wang:21arxivshort}, GLaM ~\citep{Du:21glamshort}, Chinese M6~\citep{Lin:21}, multilingual AlexaTM 20B~\citep{Soltan:22}, OPT~\citep{Zhang:22opt}, Chinchilla~\citep{Hoffmann:22short}, BLOOM~\citep{Scao:22arxivshort}, GLM-130B~\citep{Zeng:22arxivshort}, LaMDA~\citep{Thoppilan:22short}, PaLM~\citep{Chowdhery:22short}, Llama~\citep{Touvron:23arxiv}, and Gemini~\citep{GeminiTeam:23arxiv,Reid:24arxiv}.

\section{Experiments}
\label{sec:Exp}

This section presents an empirical analysis of xLSTM, emphasizing its performance in language modeling relative to established approaches. Section~\ref{sec:ExpSyntethic} focuses on elucidating xLSTM’s unique properties through evaluations on controlled synthetic tasks. In Section~\ref{sec:ExpComparison}, we report a comparative study of validation set perplexities for several state-of-the-art language modeling architectures, all trained with 15B tokens from the SlimPajama corpus~\citep{Soboleva:23}. We also carry out ablation experiments for xLSTM using this dataset. Additionally, we analyze and compare scaling trends across methods, following the methodologies set out by \citet{Kaplan:20} and \citet{Brown:20short}.

Section~\ref{sec:ExpLanguage} presents an expanded investigation in which xLSTM and the leading models from Section~\ref{sec:ExpComparison} are further trained on 300B tokens sourced from SlimPajama~\citep{Soboleva:23}. The experimental protocol includes: (1) evaluating each method’s generalization to extended sequence lengths; (2) benchmarking performance based on validation perplexity as well as a set of downstream tasks~\citep{Sutawika:24}; (3) analyzing results across the 571 distinct textual domains included in the PALOMA language benchmark suite~\citep{Magnusson:23arxivshort}; and (4) re-examining scaling properties, this time under a regime with 20-fold greater training data.

Throughout our experiments, we denote xLSTM[$a$:$b$] to represent the proportion $a/b$ between mLSTM-based and sLSTM-based xLSTM blocks. For instance, xLSTM[7:1] signifies that among every eight blocks, seven utilize mLSTM while one employs sLSTM. When the total number of blocks is set to 48, this corresponds to 42 blocks with mLSTM and 6 with sLSTM. Additionally, in all experimental configurations, mLSTM and sLSTM layers are each accompanied by up-projection blocks before and after their respective modules.

\subsection{Synthetic Tasks and Long Range Arena}
\label{sec:ExpSyntethic}
We begin by evaluating the impact of xLSTM's novel exponential gating combined with memory mixing on formal languages~\citep{Deletang:23}. Subsequently, we investigate how the new matrix memory module in xLSTM performs on the Multi-Query Associative Recall benchmark~\citep{Arora:23arxiv}. Lastly, we measure xLSTM's capability in handling extended sequence inputs using the Long Range Arena framework~\citep{Tay:21}.

\paragraph{Evaluation of xLSTM's Exponential Gating Combined with Memory Mixing.} We investigate the effectiveness of xLSTM's novel exponential gating mechanism integrated with memory mixing, a design intended to enhance performance on state tracking challenges~\citep{Merrill:24, Merrill:23}. To facilitate the assessment of extrapolation to longer input sequences, we adapt and expand the suite of formal language benchmarks originally presented by \citet{Deletang:23}. A comprehensive account of each task and a complete set of results are provided in Appendix~\ref{sec:appExpSynthetic-formalLang}. Our analysis includes a comparative study of xLSTM with a variety of architectures, such as Transformers, State Space Models, and traditional Recurrent Neural Networks. For each method, we measure the accuracy on tokens that are pertinent to the target problem, normalizing accuracy scores from 0 (random guessing) to 1 (perfect performance). The 2-block configurations evaluated comprise: xLSTM[0:1] (representing pure sLSTM), xLSTM[1:0] (pure mLSTM), xLSTM[1:1], as well as Llama, Mamba, RWKV, Retention, Hyena, standard LSTM, and LSTM layers embedded in Transformer architectures (LSTM (Block)). Figure~\ref{fig:formal-main} presents a summary of the empirical findings. Models lacking memory mixing—such as vanilla Transformers and State Space Models—consistently fail to resolve regular grammar tasks like the parity benchmark, underscoring previous observations that, in the absence of explicit state tracking, Transformers and State Space Models possess intrinsically weaker computational capabilities compared to RNNs~\citep{Merrill:24, Merrill:23, Deletang:23}.

\paragraph{Test of xLSTM's Memory Capacities on Associative Recall Tasks.} In this study, we evaluate the memory capabilities of xLSTM’s novel matrix memory using the Multi-Query Associative Recall task~\citep{Arora:23arxiv}. In each instance, a sequence is constructed with randomly selected key-value pairs from an extensive vocabulary; these pairs must be stored and subsequently retrieved. To further challenge the task compared to prior work, we scale up the number of key-value pairs to as many as 256, and lengthen the context window to a maximum of 2048. This allows for a more comprehensive assessment of the memory strengths of the competing models. Our comparison includes two-block versions of Llama, Mamba, RWKV-5, RWKV-6, xLSTM[1:1], and xLSTM[1:0]. We assess their performance based on the accuracy of recalling the associated pairs. Given that Transformers (such as Llama) possess memory that scales exponentially with the dimension of the code~\citep{Ramsauer:21}, they represent the performance benchmark for this task. The results, summarized in Figure~\ref{fig:mqar-main}, demonstrate that xLSTM[1:1] surpasses all other non-Transformer baselines, regardless of model size. Notably, including the sLSTM block does not reduce memory performance, and in fact augments it, as particularly observed in the most challenging scenario involving 256 key-value pairs. Further findings can be found in Appendix~\ref{sec:appExpSynthetic-mqar}, where additional analyses show that the increased memory of xLSTM supports extrapolation to context lengths exceeding those present during training.

\paragraph{Test of xLSTM's Long Context Capabilities on Long Range Arena.} To evaluate xLSTM's ability to process extended sequences and accommodate extensive contexts, we benchmark various approaches using the Long Range Arena~\citep{Tay:21}. Across every task, xLSTM consistently achieves high scores, indicating its architectural efficacy in managing the challenges posed by long context scenarios.

Additional information can be found in Appendix~\ref{sec:appExpSynthetic-lra}.

\subsection{Method Comparison and Ablation Study}
\label{sec:ExpComparison}

In order to investigate the central inquiry of this work—namely, the performance of our novel LSTM variants at scale in the context of language modelling—we conduct experiments in which xLSTMs, Transformers, State Space Models, and additional approaches are all trained on 15B tokens drawn from SlimPajama, using an identical auto-regressive setup. We evaluate these models using the validation set and carry out ablation experiments specifically on the xLSTM architectures.

\paragraph{Comparison of xLSTM with Existing Approaches.} To facilitate a fair comparison, we trained each model on a dataset containing 15B tokens derived from SlimPajama~\citep{Soboleva:23}, assessing their performance via perplexity on the corresponding validation set. The evaluation encompasses several techniques: xLSTM (proposed herein), GPT-3 (Transformer architecture)~\citep{Brown:20short}, Llama (Transformer)~\citep{Touvron:23arxiv}, H3 (State Space Model, SSM)~\citep{Fu:23}, Mamba (SSM)~\citep{Gu:24arxiv}, RWKV-4/5/6 (RNN family)~\citep{Peng:23arxivshort,Peng:24arxivshort}, GLA (linear Transformer)~\citep{Yang:23arxiv}, HGRN/HGRN2 (RNNs)~\citep{Qin:23,Qin:24arxiv}, RetNet (linear Transformer)~\citep{Sun:23arxiv}, Hyena (linear Transformer)~\citep{Poli:23}, and our notational variants xLSTM[1:0] and xLSTM[7:1] (as described in Section~\ref{sec:xlstm_blocks_notation}). Training was conducted using mixed precision, with the exception that models RWKV-5, RWKV-6, GLA, and HGRN2 did not employ PyTorch's automatic mixed-precision framework (refer to Appendix Section~\ref{sec:appGenTrainProc} for details). 

We classify the approaches into (a) Transformer-based models, (b) State Space Models (SSMs), and (c) both Recurrent Neural Networks (RNNs) and linear Transformers—the latter characterized by replacing the standard attention mechanism with linear operations. All models were configured to be parameter-matched to a 350M-parameter GPT-3 (using an embedding dimension of 1024 and 24 residual layers). Notably, only GPT-3 employs shared weights for input token and output embeddings, thereby reducing its parameter count. Results summarized in Table~\ref{tab:spaj15b_model_results} demonstrate that xLSTM achieves superior validation perplexity relative to all competitors. For extended discussion, please consult Appendix~\ref{sec:appExpComparison}. Figure~\ref{fig:temp_sclaw15B} illustrates the scaling characteristics observed in this experiment, suggesting that xLSTM maintains its advantage as model size increases.

\paragraph{Ablation Studies.}
As shown in Table~\ref{tab:spaj15b_model_results} and Figure~\ref{fig:temp_sclaw15B}, xLSTM delivers strong language modeling performance when trained with 15B tokens from the SlimPajama dataset. To systematically analyze the impact of each modification from LSTM to xLSTM, we incrementally convert a standard LSTM architecture into the full xLSTM design. The process begins with embedding LSTM layers within pre-LayerNorm residual frameworks. In the next phase, we modify the architecture to include a post up-projection module. The final transformation introduces exponential gating alongside the matrix memory mechanism.

Table~\ref{tab:ablstudies} (top) presents the corresponding results.

The ablation experiments demonstrate that the notable performance gains can be credited to the combination of exponential gating and the matrix memory components. Recognizing the crucial role of gating in RNNs and State Space Models, we further investigate several gating strategies. As depicted in Table~\ref{tab:ablstudies} (bottom), making each gate both learnable and input-dependent consistently enhances performance. Further analyses examining the specific backbone modules are provided in Appendix~\ref{sec:appAblDetails}.

\subsection{xLSTM as Large Language Model}
\label{sec:ExpLanguage}

To conclude our investigation, we conduct extensive language modeling experiments at scale, evaluating xLSTM's capability as a large language model (LLM). To this end, we substantially expand the training corpus, utilizing 300B tokens sourced from SlimPajama. Notably, this token budget aligns with those employed in other works, such as Mamba~\citep{Gu:24arxiv} and Griffin~\citep{De:24arxiv}.

We benchmark xLSTM against RWKV-4, Llama, and Mamba, each representing a distinct model class as described in Section~\ref{sec:ExpComparison}. RWKV-4 serves as the RNN benchmark because appropriate training precision configurations for RWKV-5, RWKV-6, and HGRN2 were identified only after commencing the 300B token-scale training (refer to Appendix~\ref{sec:appGenTrainProc}). Our experiments span multiple model scales (125M, 350M, 760M, 1.3B parameters), where we systematically evaluate each model’s length generalization, validation set accuracy, performance on a variety of downstream tasks, and language modeling ability across the 571 domains provided by the PALOMA benchmark. Additionally, we analyze their empirical scaling laws.

\paragraph{Sequence Length Extrapolation.}
\label{sec:spaj300B_ctx_extrapolation}
To begin, we evaluate how well the 1.3B-parameter versions of xLSTM, RWKV-4, Llama, and Mamba generalize to longer sequence lengths. Each model undergoes training with a context window of 2048, and their performance is subsequently assessed at extended context lengths, reaching up to 16384 tokens. Results are illustrated in Figure~\ref{fig:spaj300B_seq_extrapolation}. Notably, xLSTM models consistently achieve lower perplexity values across increased context lengths compared to the other architectures.

\paragraph{Validation Perplexity and Downstream Tasks.}
\label{sec:spaj_downstream_eval}
Next, across every model size, we assess the capabilities of xLSTM, RWKV-4, Llama, and Mamba on the SlimPajama validation dataset in terms of next token prediction, as well as on downstream evaluations that test common sense reasoning.

The third column of Table~\ref{tab:lmeval} presents the perplexity scores on the validation set achieved by various approaches. For all examined model sizes, xLSTM[1:0] and xLSTM[7:1] consistently yield the lowest validation set perplexities, indicating their superior performance. The remaining columns in Table~\ref{tab:lmeval} summarize outcomes on downstream evaluation benchmarks. xLSTM outperforms alternative methods in nearly every task and for all model scales, with the sole exception occurring on the ARC task, where Mamba occasionally attains the top result. Further details can be found in Appendix~\ref{sec:appExpLanguage}.

\paragraph{Performance on PALOMA Language Tasks.} Third, we evaluate xLSTM, RWKV-4, Llama, and Mamba models—across all parameter scales—by assessing their next token prediction accuracy on the PALOMA language tasks~\citep{Magnusson:23arxivshort}. The evaluation metric is perplexity, computed for next token prediction across 571 distinct text domains, encompassing sources from nytimes.com to Reddit's r/depression. Table~\ref{tab:ppleval} presents perplexity results, categorizing them into general language modeling domains (first seven columns) and a selection of fine-grained domain benchmarks (last five columns). Notably, xLSTM[1:0] achieves lower perplexity than xLSTM[7:1] on these language-oriented tasks.

In 568 of the 571 text domains (99.5\%), xLSTM[1:0] achieves a lower perplexity compared to Mamba. It outperforms Llama in 486 of the 571 domains (85.1\%), and surpasses RWKV-4 in 570 of 571 domains (99.8\%) with respect to perplexity. Further information is provided in Appendix~\ref{sec:appExpLanguage}.

\paragraph{Scaling Laws.} Next, we investigate the scaling properties governed by power laws, which permit the prediction of model performance at larger parameter counts~\citep{Kaplan:20,Brown:20short}. Figure~\ref{fig:temp_sclaw300B} depicts these scaling trends. Although all evaluated architectures exhibit comparable scaling patterns, their baselines differ. Among the models, RWKV-4 achieves the lowest performance, trailed by Llama and then Mamba. The xLSTM model consistently surpasses Mamba, maintaining an advantage over Mamba that mirrors Mamba’s lead over Llama. These observed scaling trends suggest that as model capacity increases, xLSTM is expected to remain superior in performance relative to both Transformer-based and State-Space approaches.

\textbf{Generation Times and Maximal Throughput.} To conclude, we evaluate the time required for text generation in Figure~\ref{fig:inference_time_generation} as well as the peak throughput presented in Figure~\ref{fig:inference_time_decoding_throughput} (left) for xLSTM variants with a parameter scale of 1.3B. The performance of these models is contrasted with implementations of Mamba, Llama, and RWKV of comparable size from HuggingFace, with the Llama model utilizing a static key-value cache. At the point of conducting these studies, full compilation of the Transformer Model's cache and employing \texttt{torch.compile} for Mamba were not functional. All models in the text generation benchmarks use a batch size of 1 and a pre-fill of 16, creating an especially advantageous condition for the Transformer model.

As depicted in Figure~\ref{fig:inference_time_generation}, the xLSTM and other recurrent architectures, namely Mamba and RWKV-4, exhibit linear scaling in generation time, whereas Llama demonstrates a quadratic scaling pattern. For the assessment of decoding throughput, batch sizes and pre-fill values are varied specifically for the Llama model.

Figure~\ref{fig:inference_time_decoding_throughput} (right) indicates that xLSTM can handle much larger batch sizes than Llama, attributed to its fixed memory requirements, which ultimately results in superior throughput relative to the alternatives.

\section{Limitations}

(i) Unlike mLSTM, sLSTM's approach to memory mixing inherently restricts parallel execution, making efficient parallel implementations unfeasible. Despite this, we have designed a custom CUDA kernel for sLSTM that achieves execution speeds within a factor of two compared to our parallelized mLSTM version. (ii) The current CUDA kernels for mLSTM lack aggressive optimization, resulting in approximately fourfold slower performance relative to either FlashAttention or the scan operation in Mamba. We expect that, with further engineering, CUDA kernels modeled after FlashAttention could significantly reduce this performance gap. (iii) Processing the $d \times d$ matrices in mLSTM's matrix memory entails considerable computational overhead. However, as the memory update and read operations are parameter-free and amenable to parallelization via standard matrix computations, the resultant increase in wall-clock time from handling complex memory structures remains modest. (iv) It is crucial to appropriately set the initial values for the forget gates. (v) Since matrix memory in mLSTM does not scale with sequence length, using much longer sequences could potentially exhaust the available memory, especially for large context windows. Nevertheless, empirical results indicate this is not a limitation for sequence contexts up to 16k tokens (see Section~\ref{sec:spaj300B_ctx_extrapolation}). (vi) Given the high computational costs of running large-scale language modeling experiments, we did not extensively optimize either the network architecture or hyperparameter configurations for the larger xLSTM models. We believe comprehensive optimization will be essential for xLSTM to achieve its optimal performance.

\section{Conclusion}
Our investigation partially addresses the straightforward inquiry: To what extent can language modeling be advanced by scaling LSTM architectures to billions of parameters? Our findings suggest that such scaled LSTMs achieve at least parity with established approaches, including Transformers and State Space Models. We introduced the xLSTM variant, incorporating exponential gating mechanisms, memory mixing, and an updated memory architecture. Empirically, xLSTM demonstrates competitive performance in language modeling relative to leading models such as Transformers and State Space Models. Scaling trends further imply that even larger xLSTM models may emerge as strong contenders against contemporary Large Language Models founded on Transformer designs. Beyond language modeling, xLSTM holds promise for applications across domains such as Reinforcement Learning, Time Series Forecasting, and the modeling of physical phenomena.

\end{document}